\section{Use one shared parse state}\label{expression-helper-state}
The private helpers share five saved registers. They are part of one internal protocol,
not independently callable public functions with the outer calling convention. The public
wrapper preserves the original caller values before assigning these parser roles.

\begin{table}[htbp]
  \centering
  \caption{Shared registers during one parse. The base and length stay fixed; the cursor and first-error fields describe progress.}
  \tagpdfsetup{table/header-rows={1}}
  \begin{tabular}{lll}
    \toprule
    Register    & Meaning                      & Helper changes          \\
    \midrule
    \texttt{s0} & Full input base              & None                    \\
    \texttt{s1} & Cursor byte offset           & Consume accepted syntax \\
    \texttt{s2} & Input length                 & None                    \\
    \texttt{s3} & Error status, initially zero & First failure           \\
    \texttt{s4} & Error byte offset            & With first failure      \\
    \bottomrule
  \end{tabular}
\end{table}

A helper returns its value in \texttt{v0} when \texttt{s3} is zero. On an error,
its caller first observes \texttt{s3} and propagates the existing diagnostic. The
public wrapper finally translates this internal protocol into the two-register
public result. In particular, a successful helper's value register is not a public
status code. Interpret registers according to the call boundary being inspected.

The scaffold gives every unfinished region a temporary branch to that helper's
balanced return. Its first unfinished parser region also supplies provisional zero.
Replace a region's placeholder rather than appending after its unconditional branch.
A partial routine may stop before consuming its complete grammar rule; the expected
observations below explicitly identify such stages. Final behavioral grading applies
only after the whole function is complete.
